%% ch06-monitoring.tex — JA4proxy Operator's User Guide
%% Chapter 6: Monitoring and Alerting

\chapter{Monitoring and Alerting}
\label{ch:monitoring}
\index{monitoring}
\index{observability}

\newthought{The observability stack is Prometheus, Grafana, Loki, and Alertmanager.}
They run in the same Compose stack as the proxy and require no external setup.
This chapter explains what each panel in the Grafana dashboard means, which
Prometheus metrics matter most, how to query logs in Loki, and what alerts are
pre-configured.

\section{The Observability Stack}
\index{Prometheus}
\index{Grafana}
\index{Loki}
\index{Alertmanager}

\begin{table}[ht]
  \small
  \begin{tabularx}{\linewidth}{llX}
    \toprule
    \textbf{Service} & \textbf{URL} & \textbf{Purpose} \\
    \midrule
    Grafana & \url{http://127.0.0.1:3000} & Dashboard, log exploration, alerting UI \\
    Prometheus & \url{http://127.0.0.1:9091} & Metrics storage and query (PromQL) \\
    Loki & \url{http://127.0.0.1:3100} & Log aggregation (queried via Grafana) \\
    Alertmanager & \url{http://127.0.0.1:9093} & Alert routing and silencing \\
    \bottomrule
  \end{tabularx}
  \caption{Observability service endpoints}
\end{table}

The proxy exposes its Prometheus metrics endpoint at \code{http://127.0.0.1:9090}.
Prometheus scrapes this every 15 seconds.\sidenote{If you run multiple proxy
instances, each instance has its own metrics endpoint. Prometheus is pre-configured
to scrape all of them via service discovery.}

\section{The Grafana Dashboard}
\index{Grafana!dashboard panels}

Open \url{http://127.0.0.1:3000} and select the \emph{JA4proxy Overview} dashboard.
The default time range is the last hour; use the time picker in the top right to
adjust.

\subsection{Connection Rate Panel}
\index{connection rate}

The top-line panel. Shows connections per second, stacked by action type.

\textbf{What to look for:}
\begin{itemize}
  \item A sudden spike in the red (BLOCKED) area indicates an active attack
    has been detected and is being blocked.
  \item A spike in total connections without a corresponding increase in BLOCKED
    connections may indicate a new attack type that the proxy is not yet classifying
    correctly — or that the dial is at zero.
  \item A sudden drop to zero in all lines means the proxy has stopped accepting
    connections. Check the container status.
\end{itemize}

\subsection{Action Distribution Panel}
\index{action distribution}

Shows the percentage breakdown of actions over the selected time window.
In a well-tuned deployment with the dial at 100, typical values are:

\begin{itemize}
  \item 60--80\% ALLOWED (legitimate traffic)
  \item 15--35\% BLOCKED (confirmed bad traffic)
  \item 1--5\% TARPITTED (high-risk, being slowed)
  \item $<$1\% RATE\_LIMITED (automated clients hitting limits)
  \item 0\% at dial=0 (everything is ALLOWED)
\end{itemize}

If BLOCKED is above 50\%, either your traffic is heavily automated or a bypass
condition needs tuning. Investigate before raising the dial further.

\subsection{Risk Score Distribution Panel}
\index{risk score!distribution}

A histogram of risk scores (0--100) across all scored connections. The shape
tells you about the quality of your traffic classification:

\begin{description}[leftmargin=3.5cm]
  \item[Bimodal] Peaks near 0--10 and 80--100. Ideal. Legitimate traffic is
    scoring low, bad traffic is scoring high, with a clear separation.
  \item[Left-skewed] Most traffic near 0, long tail to the right. Typical for
    low-attack-volume deployments. The tail represents the bad traffic.
  \item[Flat/uniform] Signals may be under-tuned or traffic is very mixed.
    Review which signals are contributing most.
  \item[Right-skewed] Most traffic scoring high. Either the proxy is seeing a
    lot of automated traffic, or whitelist entries are needed for legitimate tools.
\end{description}

\subsection{Top JA4 Fingerprints Panel}

A ranked table of fingerprints seen in the selected window, with:
\begin{itemize}
  \item Connection count
  \item Average risk score
  \item Human-readable label (if configured)
  \item Action distribution for that fingerprint
\end{itemize}

\textbf{Workflow:} Look for high-volume fingerprints with high average scores
that are currently being ALLOWED (dial=0 or dial too low). These are your
blacklist candidates. Look for high-volume fingerprints with high average scores
that are BLOCKED — verify these are genuinely malicious before keeping them blocked.

\subsection{Geographic Distribution Panel}

Connection count by country. Useful for:
\begin{itemize}
  \item Verifying your GeoIP database is working (you should see your own country)
  \item Detecting a volumetric attack from a specific region
  \item Validating country whitelist/blacklist configuration
\end{itemize}

\subsection{Beaconing Activity Panel}

Shows IPs currently flagged by the beaconing detector. Regular-interval connections
from the same IP+JA4 pair suggest C2 traffic. The panel shows the IAT coefficient
of variation — values below the threshold indicate suspicious regularity.

\subsection{Rate Limiting Panel}

Breakdown of rate-limit triggers by strategy (by\_ip, by\_ja4, by\_ip\_ja4\_pair).
High sustained activity in by\_ip suggests volumetric attacks from single sources.
High activity in by\_ja4 suggests a distributed botnet using consistent tooling.

\section{Key Prometheus Metrics}
\index{Prometheus!metrics}
\index{metrics}

\subsection{Connection Metrics}

\begin{lstlisting}[language=bash, numbers=none]
# Total connections by action
ja4proxy_connections_total{action="allowed|blocked|tarpitted|rate_limited"}

# Active connections right now
ja4proxy_active_connections

# Connection processing duration
ja4proxy_connection_duration_seconds{quantile="0.5|0.95|0.99"}
\end{lstlisting}

\subsection{Risk Scoring}

\begin{lstlisting}[language=bash, numbers=none]
# Score distribution histogram
ja4proxy_risk_score_distribution{le="10|25|50|75|100"}

# Current dial setting
ja4proxy_dial_setting
\end{lstlisting}

\subsection{Signal Performance}

\begin{lstlisting}[language=bash, numbers=none]
# Per-signal hit rates
ja4proxy_signal_hits_total{signal="abuseipdb|beaconing|asn|tls_version|..."}

# External service latency (AbuseIPDB, RDAP)
ja4proxy_external_lookup_duration_seconds{service="abuseipdb|rdap"}

# External service errors (should be near zero)
ja4proxy_external_lookup_errors_total{service="abuseipdb|rdap"}
\end{lstlisting}

\subsection{Redis Performance}

\begin{lstlisting}[language=bash, numbers=none]
# Redis operation latency
ja4proxy_redis_operation_duration_seconds{operation="get|set|zadd|..."}

# Redis errors (critical — should be zero in normal operation)
ja4proxy_redis_errors_total{operation="..."}
\end{lstlisting}

\subsection{Useful PromQL Queries}

Block rate over the last 5 minutes:
\begin{lstlisting}[language=bash, numbers=none]
rate(ja4proxy_connections_total{action="blocked"}[5m]) /
  rate(ja4proxy_connections_total[5m])
\end{lstlisting}

Average risk score (approximated from histogram):
\begin{lstlisting}[language=bash, numbers=none]
histogram_quantile(0.50,
  rate(ja4proxy_risk_score_distribution_bucket[5m]))
\end{lstlisting}

AbuseIPDB error rate (should alert if non-zero):
\begin{lstlisting}[language=bash, numbers=none]
rate(ja4proxy_external_lookup_errors_total{service="abuseipdb"}[5m])
\end{lstlisting}

\section{Loki Log Queries}
\index{Loki}
\index{LogQL}

Access Loki through the Grafana Explore interface: navigate to Explore and select
the Loki data source.

\subsection{Common LogQL Queries}

Find all blocked connections in the last hour:
\begin{lstlisting}[language=bash, numbers=none]
{container="ja4proxy"} |= "BLOCKED:"
\end{lstlisting}

Find all connections from a specific IP:
\begin{lstlisting}[language=bash, numbers=none]
{container="ja4proxy"} |= "185.220.0.1"
\end{lstlisting}

Find all connections with a specific fingerprint:
\begin{lstlisting}[language=bash, numbers=none]
{container="ja4proxy"} |= "t13d0912_a1b2c3d4e5f6"
\end{lstlisting}

Find high-scoring connections (score above 70):
\begin{lstlisting}[language=bash, numbers=none]
{container="ja4proxy"} |= "Score:" | regexp "Score: (?P<score>[789][0-9]|100)"
\end{lstlisting}

Count block events by country over time:
\begin{lstlisting}[language=bash, numbers=none]
sum by (country) (
  count_over_time(
    {container="ja4proxy"} |= "BLOCKED:" | regexp "Country: (?P<country>[A-Z]{2})"
    [5m]
  )
)
\end{lstlisting}

\section{Pre-Configured Alerts}
\index{Alertmanager}
\index{alerts}

The following alerts are pre-configured in Alertmanager. Review and adjust
thresholds to match your environment before deploying to production.

\begin{table}[ht]
  \small
  \begin{tabularx}{\linewidth}{lXl}
    \toprule
    \textbf{Alert} & \textbf{Condition} & \textbf{Severity} \\
    \midrule
    \code{ProxyDown} & Proxy metrics endpoint unreachable for 1m & Critical \\
    \code{RedisDown} & Redis health check failing for 1m & Critical \\
    \code{HighBlockRate} & Block rate $>$80\% for 5m & Warning \\
    \code{HighRiskScores} & P95 risk score $>$75 for 10m & Warning \\
    \code{AbuseIPDBErrors} & AbuseIPDB error rate $>$0 for 5m & Info \\
    \code{ScoreDrift} & Mean risk score drifts $>$20 points from 24h baseline & Warning \\
    \code{BeaconingSpike} & Beaconing detections $>$ 10× baseline for 5m & Warning \\
    \code{RedisLatency} & P99 Redis operation $>$5ms for 5m & Warning \\
    \bottomrule
  \end{tabularx}
  \caption{Pre-configured alert rules}
\end{table}

\begin{notebox}[Score drift alerting]
The \code{ScoreDrift} alert is particularly useful. A sudden increase in mean risk
scores across all connections often indicates a new attack campaign has started.
A sudden decrease may indicate a bypass condition was incorrectly configured or
that attackers have adapted their fingerprints.
\end{notebox}

\subsection{Configuring Alert Destinations}

Alert destinations (email, PagerDuty, Slack) are configured in
\code{config/alertmanager.yml}. The file ships with a commented-out template
for each integration type. \emph{See the Alertmanager documentation for
receiver configuration.}

\section{Accessing Logs Directly}

For quick investigations without the Grafana interface:

\begin{bashcode}
# All proxy logs, live
docker compose logs -f ja4proxy

# Last 500 lines, grep for errors
docker compose logs --tail=500 ja4proxy | grep -E "ERROR|WARN"

# Analytics node logs (campaign detection, slow-scan)
docker compose logs -f analytics

# HAProxy access log (raw TCP connections)
docker compose logs -f haproxy
\end{bashcode}
